% GENERATED FILE. Edit canonical lessons, then run npm run generate:books.
% canonical-source-hash: fnv1a64:572fbf8d4e00897c
% canonical-lessons: GU-C261-patang, GU-C261-udadvu, GU-C261-melo, GU-C261-garbo, GU-C261-phatakdo

\chapter{Kites and fairs}
\label{ch:gu-kites-and-fairs}
\hlchaptermodality{261}

\begin{chapteropening}
\textbf{By the end of this chapter:} \emph{I can say how we flew kites and what we did at the fair.}

Complete the last lesson of chapter 261: I can say how we flew kites and what we did at the fair.
\end{chapteropening}

\section[\texorpdfstring{\gu{પતંગ} (patang)}{પતંગ (patang)}]{\gu{પતંગ} (patang) — kite}

\label{lesson:GU-C261-patang}



\begin{quote}
Before the new one: say the Gujarati for to guess, then the Gujarati for to choose.
\end{quote}

\subsection*{You'll want to know: \gu{પતંગ}}
\textbf{\gu{પતંગ}} — \emph{patang} — \textquotedblleft{}kite\textquotedblright{}.

\begin{quote}
\textbf{\gu{મારો} \gu{પતંગ} \gu{દૂર} \gu{ગયો}.} — \emph{māro patang dūr gayo} — \textquotedblleft{}My kite went far.\textquotedblright{}
\end{quote}

\begin{cousinweb}
\textbf{\gu{પતંગ}} (\emph{patang}) comes from Sanskrit \emph{pataṅga}, \textquotedblleft{}bird; flying insect\textquotedblright{}, taken over in a learned shape and used for the paper kite (CDIAL 7721). CDIAL lists no Gujarati form.
\end{cousinweb}

\begin{grammarlens}[title={What you've built}]
One more word for everyday life at A2.
\end{grammarlens}

\subsection*{Your turn}
\begin{itemize}
\raggedright
  \item \emph{Say it:} \emph{patang}
  \item \emph{Say it:} \emph{patang}, once more
  \item \emph{Say it:} \emph{māro patang dūr gayo} — My kite went far.
  \item \emph{Recall:} say the Gujarati for to guess, then the Gujarati for to choose, then say \emph{patang} again
\end{itemize}

\begin{tcolorbox}[breakable,colback=teal!4,colframe=teal!35!black,title={Before you move on}]
What does \gu{પતંગ} mean? (\textquotedblleft{}Kite\textquotedblright{}.) Say it once more.
\end{tcolorbox}
\section[\texorpdfstring{\gu{ઉડાડવું} (uḍāḍvũ)}{ઉડાડવું (uḍāḍvũ)}]{\gu{ઉડાડવું} (uḍāḍvũ) — to fly a kite, to make fly}

\label{lesson:GU-C261-udadvu}



\begin{quote}
Before the new one: say the Gujarati for to choose, then the Gujarati for kite.
\end{quote}

\subsection*{You'll want to know: \gu{ઉડાડવું}}
\textbf{\gu{ઉડાડવું}} — \emph{uḍāḍvũ} — \textquotedblleft{}to fly a kite, to make fly\textquotedblright{}.

\begin{quote}
\textbf{\gu{અમે} \gu{બપોરે} \gu{પતંગ} \gu{ઉડાડ્યા}.} — \emph{ame bapore patang uḍāḍyā} — \textquotedblleft{}We flew kites in the afternoon.\textquotedblright{}
\end{quote}

The causative of \textbf{\gu{ઊડવું}}, \emph{ūḍvũ}, to fly up: to make something fly.

\begin{cousinweb}
\textbf{\gu{ઉડાડવું}} (\emph{uḍāḍvũ}), to make something fly, is a causative of the inherited \emph{ūḍvũ}, \textquotedblleft{}to fly up\textquotedblright{}. CDIAL lists that verb, as \emph{uḍvũ}, under Sanskrit \emph{uḍḍīyatē}, \textquotedblleft{}flies up\textquotedblright{}, and \emph{uḍāḍvũ} under its causative \emph{uḍḍāpayati} (CDIAL 1697).
\end{cousinweb}

\begin{grammarlens}[title={What you've built}]
One more word for everyday life at A2.
\end{grammarlens}

\subsection*{Your turn}
\begin{itemize}
\raggedright
  \item \emph{Say it:} \emph{uḍāḍvũ}
  \item \emph{Say it:} \emph{uḍāḍvũ}, once more
  \item \emph{Say it:} \emph{ame bapore patang uḍāḍyā} — We flew kites in the afternoon.
  \item \emph{Recall:} say the Gujarati for to choose, then the Gujarati for kite, then say \emph{uḍāḍvũ} again
\end{itemize}

\begin{tcolorbox}[breakable,colback=teal!4,colframe=teal!35!black,title={Before you move on}]
What does \gu{ઉડાડવું} mean? (\textquotedblleft{}To fly a kite, to make fly\textquotedblright{}.) Say it once more.
\end{tcolorbox}
\section[\texorpdfstring{\gu{મેળો} (meḷo)}{મેળો (meḷo)}]{\gu{મેળો} (meḷo) — fair}

\label{lesson:GU-C261-melo}



\begin{quote}
Before the new one: say the Gujarati for kite, then the Gujarati for to fly a kite.
\end{quote}

\subsection*{You'll want to know: \gu{મેળો}}
\textbf{\gu{મેળો}} — \emph{meḷo} — \textquotedblleft{}fair\textquotedblright{}.

\begin{quote}
\textbf{\gu{અમે} \gu{ગામના} \gu{મેળામાં} \gu{ગયા}.} — \emph{ame gāmnā meḷāmā̃ gayā} — \textquotedblleft{}We went to the village fair.\textquotedblright{}
\end{quote}

\begin{cousinweb}
\textbf{\gu{મેળો}} (\emph{meḷo}) is a tadbhava. It comes from Sanskrit \emph{mēla}, \textquotedblleft{}meeting\textquotedblright{}; CDIAL glosses the Gujarati word \textquotedblleft{}assembly, fair\textquotedblright{} (CDIAL 10331).
\end{cousinweb}

\begin{grammarlens}[title={What you've built}]
One more word for everyday life at A2.
\end{grammarlens}

\subsection*{Your turn}
\begin{itemize}
\raggedright
  \item \emph{Say it:} \emph{meḷo}
  \item \emph{Say it:} \emph{meḷo}, once more
  \item \emph{Say it:} \emph{ame gāmnā meḷāmā̃ gayā} — We went to the village fair.
  \item \emph{Recall:} say the Gujarati for kite, then the Gujarati for to fly a kite, then say \emph{meḷo} again
\end{itemize}

\begin{tcolorbox}[breakable,colback=teal!4,colframe=teal!35!black,title={Before you move on}]
What does \gu{મેળો} mean? (\textquotedblleft{}Fair\textquotedblright{}.) Say it once more.
\end{tcolorbox}
\section[\texorpdfstring{\gu{ગરબો} (garbo)}{ગરબો (garbo)}]{\gu{ગરબો} (garbo) — garba, a circle dance}

\label{lesson:GU-C261-garbo}



\begin{quote}
Before the new one: say the Gujarati for to fly a kite, then the Gujarati for fair.
\end{quote}

\subsection*{You'll want to know: \gu{ગરબો}}
\textbf{\gu{ગરબો}} — \emph{garbo} — \textquotedblleft{}garba, a circle dance\textquotedblright{}.

\begin{quote}
\textbf{\gu{અમે} \gu{રાત્રે} \gu{ગરબો} \gu{ગાયો}.} — \emph{ame rātre garbo gāyo} — \textquotedblleft{}We sang a garba at night.\textquotedblright{}
\end{quote}

\begin{cousinweb}
The origin of \textbf{\gu{ગરબો}} (\emph{garbo}) is not settled. A link with Sanskrit \emph{gárbha}, \textquotedblleft{}womb, inside\textquotedblright{}, is often proposed, but CDIAL lists no Gujarati \emph{garbo} under that word (CDIAL 4055).
\end{cousinweb}

\begin{grammarlens}[title={What you've built}]
One more word for everyday life at A2.
\end{grammarlens}

\subsection*{Your turn}
\begin{itemize}
\raggedright
  \item \emph{Say it:} \emph{garbo}
  \item \emph{Say it:} \emph{garbo}, once more
  \item \emph{Say it:} \emph{ame rātre garbo gāyo} — We sang a garba at night.
  \item \emph{Recall:} say the Gujarati for to fly a kite, then the Gujarati for fair, then say \emph{garbo} again
\end{itemize}

\begin{tcolorbox}[breakable,colback=teal!4,colframe=teal!35!black,title={Before you move on}]
What does \gu{ગરબો} mean? (\textquotedblleft{}Garba, a circle dance\textquotedblright{}.) Say it once more.
\end{tcolorbox}
\section[\texorpdfstring{\gu{ફટાકડો} (phaṭākḍo)}{ફટાકડો (phaṭākḍo)}]{\gu{ફટાકડો} (phaṭākḍo) — firecracker}

\label{lesson:GU-C261-phatakdo}



\begin{quote}
Before the new one: say the Gujarati for fair, then the Gujarati for garba.
\end{quote}

\subsection*{You'll want to know: \gu{ફટાકડો}}
\textbf{\gu{ફટાકડો}} — \emph{phaṭākḍo} — \textquotedblleft{}firecracker\textquotedblright{}.

\begin{quote}
\textbf{\gu{દિવાળીમાં} \gu{મેં} \gu{ફટાકડા} \gu{ખરીદ્યા}.} — \emph{divāḷīmā̃ mẽ phaṭākḍā kharīdyā} — \textquotedblleft{}I bought firecrackers at Diwali.\textquotedblright{}
\end{quote}

\begin{cousinweb}
The origin of \textbf{\gu{ફટાકડો}} (\emph{phaṭākḍo}) is not settled. It looks built on a sound-word for a bang, but CDIAL has no entry for it and our sources do not trace it.
\end{cousinweb}

\begin{grammarlens}[title={What you've built}]
One more word for everyday life at A2.
\end{grammarlens}

\subsection*{Your turn}
\begin{itemize}
\raggedright
  \item \emph{Say it:} \emph{phaṭākḍo}
  \item \emph{Say it:} \emph{phaṭākḍo}, once more
  \item \emph{Say it:} \emph{divāḷīmā̃ mẽ phaṭākḍā kharīdyā} — I bought firecrackers at Diwali.
  \item \emph{Recall:} say the Gujarati for fair, then the Gujarati for garba, then say \emph{phaṭākḍo} again
\end{itemize}

\begin{tcolorbox}[breakable,colback=teal!4,colframe=teal!35!black,title={Before you move on}]
What does \gu{ફટાકડો} mean? (\textquotedblleft{}Firecracker\textquotedblright{}.) Say it once more.
\end{tcolorbox}
